% GENERATED FILE. Edit canonical lessons, then run npm run generate:books.
% canonical-source-hash: fnv1a64:b65b079a19ea59e0
% canonical-lessons: SA-C241-krte, SA-C241-yena, SA-C241-tena, SA-C241-yat, SA-C241-parantu

\chapter{For the sake of, By which, By that, That, However}
\label{ch:sa-for-the-sake-of-by-which-by-that-that-however}
\hlchaptermodality{241}

\begin{chapteropening}
\textbf{By the end of this chapter:} \emph{I can say for the sake of, by which, by that, that and however.}

Complete the last lesson of chapter 241: I can say for the sake of, by which, by that, that and however.
\end{chapteropening}

\section[\texorpdfstring{\sk{कृते} (kṛte)}{कृते (kṛte)}]{\sk{कृते} (kṛte) — for the sake of}

\label{lesson:SA-C241-krte}



\begin{quote}
Before the new one: say the Sanskrit for daily, then the Sanskrit for constantly, always.
\end{quote}

\subsection*{You'll want to know: \sk{कृते}}
\textbf{\sk{कृते}} — \emph{kṛte} — \textquotedblleft{}for the sake of\textquotedblright{}.

\begin{grammarlens}[title={What you've built}]
One more word for this chapter.
\end{grammarlens}

\subsection*{Your turn}
\begin{itemize}
\raggedright
  \item \emph{Say it:} \emph{kṛte}
  \item \emph{Say it:} \emph{kṛte}, once more
  \item \emph{Say it:} say \emph{kṛte}
  \item \emph{Recall:} say the Sanskrit for daily, then the Sanskrit for constantly, always, then say \emph{kṛte} again
\end{itemize}

\begin{tcolorbox}[breakable,colback=teal!4,colframe=teal!35!black,title={Before you move on}]
What does \sk{कृते} mean? (\textquotedblleft{}For the sake of\textquotedblright{}.) Say it once more.
\end{tcolorbox}
\section[\texorpdfstring{\sk{येन} (yena)}{येन (yena)}]{\sk{येन} (yena) — by which, so that}

\label{lesson:SA-C241-yena}



\begin{quote}
Before the new one: say the Sanskrit for constantly, always, then the Sanskrit for for the sake of.
\end{quote}

\subsection*{You'll want to know: \sk{येन}}
\textbf{\sk{येन}} — \emph{yena} — \textquotedblleft{}by which, so that\textquotedblright{}.

\begin{grammarlens}[title={What you've built}]
One more word for this chapter.
\end{grammarlens}

\subsection*{Your turn}
\begin{itemize}
\raggedright
  \item \emph{Say it:} \emph{yena}
  \item \emph{Say it:} \emph{yena}, once more
  \item \emph{Say it:} say \emph{yena}
  \item \emph{Recall:} say the Sanskrit for constantly, always, then the Sanskrit for for the sake of, then say \emph{yena} again
\end{itemize}

\begin{tcolorbox}[breakable,colback=teal!4,colframe=teal!35!black,title={Before you move on}]
What does \sk{येन} mean? (\textquotedblleft{}By which, so that\textquotedblright{}.) Say it once more.
\end{tcolorbox}
\section[\texorpdfstring{\sk{तेन} (tena)}{तेन (tena)}]{\sk{तेन} (tena) — by that, therefore}

\label{lesson:SA-C241-tena}



\begin{quote}
Before the new one: say the Sanskrit for for the sake of, then the Sanskrit for by which.
\end{quote}

\subsection*{You'll want to know: \sk{तेन}}
\textbf{\sk{तेन}} — \emph{tena} — \textquotedblleft{}by that, therefore\textquotedblright{}.

\begin{grammarlens}[title={What you've built}]
One more word for this chapter.
\end{grammarlens}

\subsection*{Your turn}
\begin{itemize}
\raggedright
  \item \emph{Say it:} \emph{tena}
  \item \emph{Say it:} \emph{tena}, once more
  \item \emph{Say it:} say \emph{tena}
  \item \emph{Recall:} say the Sanskrit for for the sake of, then the Sanskrit for by which, then say \emph{tena} again
\end{itemize}

\begin{tcolorbox}[breakable,colback=teal!4,colframe=teal!35!black,title={Before you move on}]
What does \sk{तेन} mean? (\textquotedblleft{}By that, therefore\textquotedblright{}.) Say it once more.
\end{tcolorbox}
\section[\texorpdfstring{\sk{यत्} (yat)}{यत् (yat)}]{\sk{यत्} (yat) — that (introducing a clause)}

\label{lesson:SA-C241-yat}



\begin{quote}
Before the new one: say the Sanskrit for by which, then the Sanskrit for by that.
\end{quote}

\subsection*{You'll want to know: \sk{यत्}}
\textbf{\sk{यत्}} — \emph{yat} — \textquotedblleft{}that (introducing a clause)\textquotedblright{}.

\begin{grammarlens}[title={What you've built}]
One more word for this chapter.
\end{grammarlens}

\subsection*{Your turn}
\begin{itemize}
\raggedright
  \item \emph{Say it:} \emph{yat}
  \item \emph{Say it:} \emph{yat}, once more
  \item \emph{Say it:} say \emph{yat}
  \item \emph{Recall:} say the Sanskrit for by which, then the Sanskrit for by that, then say \emph{yat} again
\end{itemize}

\begin{tcolorbox}[breakable,colback=teal!4,colframe=teal!35!black,title={Before you move on}]
What does \sk{यत्} mean? (\textquotedblleft{}That (introducing a clause)\textquotedblright{}.) Say it once more.
\end{tcolorbox}
\section[\texorpdfstring{\sk{परन्तु} (parantu)}{परन्तु (parantu)}]{\sk{परन्तु} (parantu) — however, but}

\label{lesson:SA-C241-parantu}



\begin{quote}
Before the new one: say the Sanskrit for by that, then the Sanskrit for that.
\end{quote}

\subsection*{You'll want to know: \sk{परन्तु}}
\textbf{\sk{परन्तु}} — \emph{parantu} — \textquotedblleft{}however, but\textquotedblright{}.

\begin{grammarlens}[title={What you've built}]
One more word for this chapter.
\end{grammarlens}

\subsection*{Your turn}
\begin{itemize}
\raggedright
  \item \emph{Say it:} \emph{parantu}
  \item \emph{Say it:} \emph{parantu}, once more
  \item \emph{Say it:} say \emph{parantu}
  \item \emph{Recall:} say the Sanskrit for by that, then the Sanskrit for that, then say \emph{parantu} again
\end{itemize}

\begin{tcolorbox}[breakable,colback=teal!4,colframe=teal!35!black,title={Before you move on}]
What does \sk{परन्तु} mean? (\textquotedblleft{}However, but\textquotedblright{}.) Say it once more.
\end{tcolorbox}
